\section{The Mercy Beside the Cross}

Luke recounts a conversion in the last hours of a man's life. Two criminals are led out with Jesus and crucified on either side of him. The rulers and soldiers mock the Savior, challenging him to prove his identity by saving himself. One of the condemned joins that demand. The other turns toward Jesus with a different appeal (Lk 23:32--43). They suffer the same sentence and see the same crucified Lord, yet respond differently to him. The penitent's response gives us reason to hope for those who turn late and to consider how we are responding to grace now.

The penitent begins by addressing his fellow sufferer. He recalls the fear of God, acknowledges that their own sentence follows their deeds, and bears witness that Jesus has done nothing wrong. Luke leaves the detailed history of their offenses untold. He records the penitent's acknowledgment of responsibility and his testimony to Christ's innocence. The man answers his companion's mockery by admitting the guilt they share and distinguishing Jesus from them. His suffering has not prevented this act of repentance and faith.

We can practice this honesty before we are near our own death. We may be occupied with the unfairness of particular circumstances while avoiding the wrong for which we really are responsible. An accusation may be exaggerated; the accuser may have faults; the consequences may hurt more than we expected. Those questions may matter, yet none removes the need to acknowledge the sin itself. The thief admits that his deeds deserve punishment while bearing witness to Jesus' innocence. We too can acknowledge what we have done without waiting for everyone else to acknowledge their faults or for the consequences to become easier to bear.

He then addresses Jesus as a king, asking to be remembered when Jesus comes into his kingdom. What stands before his bodily eyes is a man sharing the shame of execution, mocked beneath a royal inscription. His appeal reaches beyond that visible humiliation. He entrusts his future to the sufferer beside him, believing that the crucified Jesus has a kingdom in which he can remember another. The request does not explain how this will be accomplished. Faith can turn toward the person of the Savior before it understands the manner in which his promise will be fulfilled.

Jesus promises that the penitent will be with him in paradise that very day (Lk 23:43). The man who faces execution will live in communion with Christ. He has little time left for action, yet can entrust himself to the Lord who receives him. His hope depends on the authority and mercy of Jesus. The mockery continues around the cross, but Jesus answers the man's appeal as the King he has confessed, promising what only the Savior can give.

We should let this scene give its full consolation to those who have turned late. A life may contain grave wrongs and long resistance. Its repair may be incomplete; the time remaining may be short. The penitent does not have to persuade the Savior that the past was better than it was. He can acknowledge it and turn. The Church's urgency in bringing Penance to the seriously ill belongs to this confidence in Christ's mercy. No human calculation of years wasted gives us authority to tell a contrite person that it is now too late to seek forgiveness.

Deliberately delaying conversion differs from the thief's response. He turns to Jesus in the opportunity given to him. His repentance is reason for us to turn too; we have no promise that we shall control the conditions of our own last hours. Admiration for his trust becomes fruitful when we act on the grace already offered. Whoever can confess, make restitution, or cease a wrong has a present opportunity to repent. We can ask Christ to help us use it, with the same confidence in his mercy that the thief expressed beside him.

When time and ability remain, a sincere return also seeks to repair the injury that can still be repaired. A dying person may be able to acknowledge a neglected child, correct a false accusation, disclose a debt, or ask pardon from someone long wounded. Such acts can be costly even within weakness. They express contrition through attention to the person who has been injured. Where possible, those helping the dying can assist with the conversation or practical arrangement required. When circumstances make a desired reparation impossible, the person can place both the sorrow and the unfulfilled duty before God, who knows what can and cannot still be done.

The Church and the family can help this turning by making room for honest speech. A familiar attempt to reassure the dying is to insist that they have never done anything requiring serious repentance. Affection may intend comfort, yet the sufferer's conscience may be asking for help to acknowledge a real fault. We can listen without shock, help obtain a priest, and allow a person to ask forgiveness without being hurried back into a flattering account of life. Christ answers the thief's acknowledgment of guilt with a promise of paradise. That promise gives us reason to help the contrite person seek forgiveness confidently.

At the same time, we cannot prescribe another person's final words. Some will be unable to speak; some will lack the strength to explain what is happening inwardly; some will die without the opportunity for a recognizable farewell. The Gospel scene directs us toward the Lord who receives the penitent, whose knowledge is greater than our observation. We offer every help within our power and pray for the grace of repentance and perseverance. The absence of a scene we can understand does not entitle us to describe the secret encounter between God and a soul.

Luke then recounts the Lord's own final surrender. After the promise to the thief, darkness falls, the veil of the temple is torn, and Jesus commends his spirit to the Father (Lk 23:44--46). The penitent has entrusted himself to Jesus; Jesus freely surrenders his human life to the Father. One is a sinner's appeal for mercy, the other the Son's saving obedience. In considering them together, we can recognize our dependence on what Christ accomplishes for us. The man who asks to be remembered receives the promise of salvation from the Lord who is giving his life to save him.

When we read this account again, we can acknowledge a particular wrong before Christ and ask for his mercy. The past may be painful to confess, and a desired reparation may be difficult. The thief's appeal shows that we can approach the Savior with that need plainly admitted. We can also pray for someone who appears far from repentance, asking that grace reach him and that we be ready to help if he seeks a way back.

